\section{Introduction}
\label{sec:introduction}

The Basis Network enterprise zkEVM L2 architecture assigns each
enterprise its own chain with an independent sequencer and prover. Each
chain produces halo2-KZG validity proofs that are verified on the Basis
Network L1 (an Avalanche Subnet-EVM instance with sub-second finality).
This per-enterprise isolation provides strong privacy guarantees: no
enterprise exposes its transaction data or state transitions to any
other participant.

However, independent verification introduces a scalability problem.
Each halo2-KZG proof verification costs approximately 420K gas on the
EVM~\cite{axiom-v2, orbiter-gas}. For $N$ enterprises, total
verification gas scales as $N \times 420\text{K}$, reaching 3.36M gas
at $N{=}8$ and 6.72M gas at $N{=}16$. Although the Basis Network L1
operates with zero gas fees, the computational cost of verification
still consumes block capacity and imposes throughput limits. More
critically, as the network scales to dozens or hundreds of enterprises,
per-block verification becomes a bottleneck regardless of fee structure.

Recursive proof composition offers a solution: aggregate $N$ independent
proofs into a single proof verifiable in constant gas, amortizing
verification cost across all participants. The zero-knowledge proof
literature provides three principal aggregation strategies:

\begin{enumerate}
  \item \textbf{Binary tree accumulation.} Recursively verify pairs of
    proofs inside a circuit, building a tree of depth
    $\lceil \log_2 N \rceil$. The final root proof attests to the
    validity of all leaves. This approach is production-proven at Scroll
    (chunk-batch-bundle hierarchy)~\cite{scroll-darwin} and Axiom
    (snark-verifier)~\cite{axiom-v2, snark-verifier}.

  \item \textbf{Folding schemes.} Fold $N$ proof instances into a single
    accumulated instance without performing in-circuit verification,
    then produce a single ``decider'' proof. Nova~\cite{nova},
    ProtoGalaxy~\cite{protogalaxy}, and HyperNova~\cite{hypernova}
    achieve this with $O(1)$ memory and proof size regardless of~$N$.

  \item \textbf{Batch verification.} Exploit algebraic structure
    (e.g., pairing batching) to verify $N$ proofs with sub-linear
    pairing operations. SnarkPack~\cite{snarkpack} achieves
    $O(\log N)$ proof size for Groth16 batches.
\end{enumerate}

This paper presents a systematic evaluation of all three strategies
applied to the Basis Network enterprise aggregation problem.

\subsection{Contributions}

\begin{itemize}
  \item A comparative benchmark of three aggregation strategies across
    $N \in \{1, 2, 4, 8, 16\}$ enterprises, measuring proof size,
    aggregation time, and L1 verification gas with 30 stochastic
    iterations per configuration.
  \item Quantification of gas savings: folding with Groth16 decider
    achieves $15.3\times$ reduction at $N{=}8$ (from 3.36M to 220K
    gas), exceeding the $N$-fold hypothesis.
  \item Identification of four aggregation invariants
    (Section~\ref{sec:invariants}) required for production soundness.
  \item A phased deployment strategy balancing gas efficiency against
    tooling maturity (Section~\ref{sec:deployment}).
  \item Reconciliation of all measurements against published benchmarks
    from six production systems with no divergence exceeding $10\times$.
\end{itemize}

\subsection{Paper Organization}

Section~\ref{sec:related} surveys recursive SNARK composition, folding
schemes, and production aggregation systems.
Section~\ref{sec:system-model} defines the enterprise aggregation model
and its constraints. Section~\ref{sec:methodology} describes the
experimental design, benchmark protocol, and anti-confirmation bias
measures. Section~\ref{sec:results} presents the complete benchmark
results. Section~\ref{sec:discussion} analyzes trade-offs, identifies
invariants, and proposes a deployment strategy.
Section~\ref{sec:conclusion} summarizes findings and outlines future
work.
